
\section{致谢}
\label{sec:1-6}

在研究一般拉格朗日子流形的弗洛尔理论时，构造一列障碍类，并在这些障碍类平凡时形变弗洛尔边界映射这一思想，得自孔采维奇。我们还从他的论文中学到了许多内容；这些内容在本书若干处起着实质性的作用。我们感谢 Urs Frauenfelder 指出了本书 2000 年版中与阿诺尔德--纪梵塔尔猜想的证明有关的一处实质性错误。我们也感谢 Cheol-Hyun Cho 所作的若干宝贵讨论，尤其是关于量子上同调与对角子流形的弗洛尔上同调之间同构的讨论；他的一条意见还促使我们写成第 7.4.2--7.4.11 小节。在本书的撰写或修订过程中，堀健太朗、斯塔谢夫、南範彦（Norihiko Minami；日文 南範彦）、古田幹雄（Mikio Furuta；日文 古田幹雄）、牛腸徹（Toru Gocho；日文 牛腸徹）、河野明（Akira Kono；日文 河野明）和加藤文元（Fumiharu Kato；日文 加藤文元）分别为我们提供了关于势函数与朗道--金兹堡模型、$A_\infty$ 代数、有理同伦论、重整化理论、非交换辛几何与 $A_\infty$ 代数的关系、环路空间同调以及刚性解析几何的各种有用信息，我们也谨向他们致谢。

我们于 1998 年开始撰写本书，初步版本在 2000 年 12 月完成。自那以后，我们一直把它放在第一署名作者的主页（\url{http://www.math.kyoto-u.ac.jp/~fukaya/}）上。

自本书 2000 年版问世以来，为完成这一最终版本而进行的修订，耗时远超我们最初的预期。在此期间，本书的许多结果在作者们的各篇文章中得到进一步阐释（见\textbf{[ChOh03, Fuk03II,05II, Oht01]}）。与 2000 年版相比，本书内容已有大量修改和改进。因此，有必要对 2000 年版与目前最终版本之间的关系作些说明。

首先需要指出，2000 年版中有若干错误。其中最明显的一处是：在撰写本书时，作者要证明阿诺尔德--纪梵塔尔猜想，还需要对 $(M,\omega)$ 附加球面正性条件。除这一附加条件外，2000 年版导言所陈述各定理的证明中的所有错误，均已在本书或此前的若干论文\textbf{[FOOO09I,09II]} 中得到纠正。我们强调，2000 年版导言中陈述的所有定理，除其中的“定理 H”需要增加一个假设外，都仍按原陈述成立（见\textbf{[FOOO09I]}）。

对于读过 2000 年版或 2006 年版、并希望了解 2000 年版与这一最终版本之间差异的读者，我们准备了说明这些差异的文献\textbf{[FOOO09III]}。事实上，由于出版者对 AMS/IP 系列篇幅的限制，先前 2007 年版所包含的两章未收入这一最终版本：其中一章在拉格朗日子流形满足额外的“球面正性”假设下，研究取 $\mathbb Z$ 或 $\mathbb Z_2$ 系数的拉格朗日弗洛尔上同调；另一章研究拉格朗日弗洛尔上同调与拉格朗日手术之间的关系。这两章将作为独立论文另行发表\textbf{[FOOO09I,09II]}。不过我们要指出，2007 年版中的这两章已经是最终定稿。从本书移去这两章带来的一个缺点，是书中没有收入很多拉格朗日弗洛尔上同调的计算实例。更多这类计算实例，参见\textbf{[FOOO08I,08II,09I,09II]}。

第一署名作者深谷贤治（K.F.）感谢 IHES；他于 1997 年访问 IHES 时，曾在那里与孔采维奇就本项目作了富有成果的讨论。第一、第三和第四署名作者（K.F.、H.O.、K.O.，即深谷贤治、太田啓史（Hiroshi Ohta）和小野薫）也感谢吴龙根以及 KIAS 在他们 1999 年和 2000 年访问 KIAS 期间给予的款待。他们还感谢 ETH、华威大学、约翰斯·霍普金斯大学以及波恩马克斯·普朗克数学研究所在访问期间给予的款待。这些机构提供了优良的工作环境，使他们得以专注于本项目；他们从这些访问中获益良多。第二署名作者吴龙根（Y.O.）感谢深谷贤治以及 RIMS 在他 1999 年秋季逗留 RIMS 期间给予的经费支持和款待。第四署名作者小野薫感谢 IHES 在他 1998 年夏季逗留期间给予的款待。

作者感谢丘成桐教授建议将本书作为 \emph{AMS/IP Studies of Mathematics} 丛书的一卷出版。作者也感谢 Tanaka 女士为本书绘制了精美的插图。
